\documentclass[
	% -- opções da classe memoir --
	article,			% indica que é um artigo acadêmico
	11pt,				% tamanho da fonte
	oneside,			% para impressão apenas no verso. Oposto a twoside
	a4paper,			% tamanho do papel. 
	% -- opções da classe abntex2 --
	%chapter=TITLE,		% títulos de capítulos convertidos em letras maiúsculas
	%section=TITLE,		% títulos de seções convertidos em letras maiúsculas
	%subsection=TITLE,	% títulos de subseções convertidos em letras maiúsculas
	%subsubsection=TITLE % títulos de subsubseções convertidos em letras maiúsculas
	% -- opções do pacote babel --
	english,			% idioma adicional para hifenização
	brazil,				% o último idioma é o principal do documento
	sumario=tradicional
	]{abntex2}

% ---
% PACOTES
% ---

% ---
% Pacotes fundamentais 
% ---
\usepackage{lmodern}			% Usa a fonte Latin Modern
\usepackage[T1]{fontenc}		% Selecao de codigos de fonte.
\usepackage[utf8]{inputenc}		% Codificacao do documento (conversão automática dos acentos)
\usepackage{indentfirst}		% Indenta o primeiro parágrafo de cada seção.
\usepackage{nomencl} 			% Lista de simbolos
\usepackage{color}				% Controle das cores
\usepackage{graphicx}			% Inclusão de gráficos
\usepackage{microtype} 			% para melhorias de justificação
\usepackage[table]{xcolor}      % Para cor das células das tabelas
\usepackage{enumitem}
\usepackage{float}              % Para o posicionamento [H] da tabela
\usepackage{longtable}          % Para tabelas longas
% ---
% Pacotes adicionais, usados apenas no âmbito do Modelo Canônico do abnteX2
% ---
\usepackage{lipsum}				% para geração de dummy text
% ---
		
% ---
% Pacotes de citações
% ---
\usepackage[brazilian,hyperpageref]{backref}	 % Paginas com as citações na bibl
\usepackage[alf]{abntex2cite}	% Citações padrão ABNT
% ---

% ---
% Configurações do pacote backref
% Usado sem a opção hyperpageref de backref
\renewcommand{\backrefpagesname}{Citado na(s) página(s):~}
% Texto padrão antes do número das páginas
\renewcommand{\backref}{}
% Define os textos da citação
\renewcommand*{\backrefalt}[4]{
	\ifcase #1 %
		Nenhuma citação no texto.%
	\or
		Citado na página #2.%
	\else
		Citado #1 vezes nas páginas #2.%
	\fi}%
% ---

% ---
% Informações de dados para CAPA e FOLHA DE ROSTO
% ---
\titulo{\textbf{Requisitos e Análise em Engenharia de Software:\\Sistema de Intermediação de Brechós}}
\autor{
    Caroline Santos, Beatriz Eduão, Ian Daniel 
	\\[0.2cm] 
	João Pedro M., Arthur Soares, Luiz Manoel
}
\local{Brasil}
\data{Maio, 2026}
% ---

% ---
% Configurações de aparência do PDF final

% alterando o aspecto da cor azul
\definecolor{blue}{RGB}{41,5,195}

% informações do PDF
\makeatletter
\hypersetup{
        %pagebackref=true,
		pdftitle={\@title}, 
		pdfauthor={\@author},
        pdfsubject={Modelo de artigo científico com abnTeX2},
        pdfcreator={LaTeX with abnTeX2},
		pdfkeywords={abnt}{latex}{abntex}{abntex2}{atigo científico}, 
		colorlinks=true,            % false: boxed links; true: colored links
        linkcolor=blue,             % color of internal links
        citecolor=blue,             % color of links to bibliography
        filecolor=magenta,          % color of file links
		urlcolor=blue,
		bookmarksdepth=4
}
\makeatother
% --- 

% ---
% compila o indice
% ---
\makeindex
% ---

% ---
% Altera as margens padrões
% ---
\setlrmarginsandblock{3cm}{3cm}{*}
\setulmarginsandblock{3cm}{3cm}{*}
\checkandfixthelayout
% ---

% --- 
% Espaçamentos entre linhas e parágrafos 
% --- 

% O tamanho do parágrafo é dado por:
\setlength{\parindent}{1.3cm}

% Controle do espaçamento entre um parágrafo e outro:
\setlength{\parskip}{0.2cm}

% Espaçamento simples
\SingleSpacing

% ----
% Início do documento
% ----
\begin{document}

\frenchspacing 

% ----------------------------------------------------------
% ELEMENTOS PRÉ-TEXTUAIS
% ----------------------------------------------------------
\newcommand{\inst}{UNIVERSIDADE FEDERAL DE SERGIPE (UFS)}
\newcommand{\centro}{CENTRO DE CIÊNCIAS EXATAS E TECNOLOGIA (CCET)}
\newcommand{\dept}{DEPARTAMENTO DE COMPUTAÇÃO (DCOMP)}
\newcommand{\disc}{ENGENHARIA DE SOFTWARE - COMP0503}

\makeatletter
\def\@maketitle{
\newpage
\null
\vskip 2em
\begin{center}
    \normalsize
    {\inst \par \centro \par \dept \par \disc \par}
    \vskip 3em
\let \footnote \thanks
{\normalsize \@title \par}
\vskip 3em
{\normalsize \flushright
    \lineskip .5em
    \begin{tabular}[t]{c}
    \@author
    \end{tabular}\par}
\vskip 1em
\end{center}%
\par
\vskip 1.5em}
\makeatother

\maketitle

\thispagestyle{empty}

\begin{resumoumacoluna}[\normalsize \textsf{\textbf{RESUMO}}]
Este documento compõe o projeto de Sistema de Intermediação de Brechós destinado à disciplina de Engenharia de Software na Universidade Federal de Sergipe (UFS), ministrada pelo Prof. Dr. Michel Soares. O objetivo principal se pauta em estabelecer uma análise sobre os requisitos de um sistema para gerenciamento de brechós no modelo C2C (\textit{Consumer-to-Consumer}), aplicando princípios da Engenharia de Requisitos. A documentação contempla estruturalmente o padrão estabelecido pela IEEE 830/1998 de práticas recomendadas para especificação de requisitos de Software.
 
 \vspace{\onelineskip}
 
 \noindent
 \textbf{Palavras-chaves}: requisitos. análise. brechós. engenharia de requisitos. engenharia de software.
\end{resumoumacoluna}

\textual

\clearpage
\setcounter{page}{1}

\section{Introdução}
\label{sec:Introducao}

\subsection{Propósito do documento}
\label{sec:Proposito}

O documento intenciona aplicar técnicas de Requisitos e Análise em projetos de Engenharia de Software sobre um Sistema de Intermediação de Brechós, no modelo C2C (Consumer-to-Consumer). Enquanto uma Especificação de Requisitos de Software (\textit{Software Requirements Specification} – SRS), a documentação em questão também visa definir de forma clara o próposito, as funcionalidades e as limitações do projeto.

Para tal, a sua estrutura fundamenta-se no padrão IEEE 830/1998, o qual estabelece diretrizes e práticas recomendadas para a representação de requisitos de software de maneira organizada, consistente e completa, além de incorporar princípios gerais da Engenharia de Requisitos. O público esperado desta SRS abrange: equipe responsável pela implementação do sistema, equipe de testes e garantia de qualidade do sistema, gerentes de projeto, clientes e Stakeholders beneficiados pela iniciativa.

Ao final deste documento, espera-se produzir uma ferramenta acurada e útil sobre as necessidades dos usuários, servindo como alicerce para as fases subsequentes do desenvolvimento do sistema proposto.

\subsection{Escopo do produto}
\label{sec:Escopo}

\subsubsection{Produtos de software utilizados}
\label{sec:Produtos}

O desenvolvimento do Sistema de Intermediação de Brechós poderá utilizar ferramentas modernas de desenvolvimento web, banco de dados e versionamento de código. Entre os produtos de software utilizados, destacam-se ambientes de desenvolvimento integrados (IDEs), sistemas gerenciadores de banco de dados relacionais ou não relacionais, bibliotecas de interface gráfica para aplicações web e plataformas de hospedagem em nuvem.

Além disso, o sistema poderá integrar serviços externos responsáveis pelo processamento de pagamentos, autenticação de usuários e armazenamento de imagens enviadas pelos anunciantes.

\subsubsection{O que o sistema realiza (e o que não realiza)}
\label{sec:Realiza}

Respeitando a sua natureza C2C, o sistema deve intermediar o processo de compra e venda de peças de vestuário entre usuários, no qual os usuários podem exercer ambos papéis de anunciantes e/ou consumidores conforme sua preferência. O sistema deve automatizar a conexão entre usuários que desejam comprar e vender roupas usadas, promovendo os brechós, tradicionalmente manuais, como uma nova experiência organizada e intuitiva. Além disso, o sistema deve possibilitar o gerenciamento de informações pessoais do usuário (nome, e-mail, senha), permitir o anúncio de novos itens ou a exclusão de itens indisponíveis e implementar uma busca personalizada por categorias de peças de vestuário (camiseta, casaco, vestido).

O sistema busca solucionar as problemáticas de: ineficiência no anúncio manual de peças de vestuário, dificuldade de descoberta de itens específicos pelos usuários, falha no canal de comunicação entre anunciantes e consumidores, inconsistência em catálagos de peças anunciadas e ausência de uma plataforma dedicada ao processo de compra e venda de roupas usadas em brechós.

Nesse sentido, o propósito do sistema não é realizar a validação de qualidade, durabilidade ou integridade das roupas. O sistema também não é responsável por garantir que a peça de roupa sirva conforme o esperado ao cliente. Entretanto, o sistema deve permitir o upload de imagens fotográficas nos anúncios dos itens, cabendo ao usuário o julgamento de quais itens o convém.

\subsubsection{Descrição da aplicação de software proposta}
\label{sec:DescricaoAplicacao}

Inicialmente, a aplicação deve funcionar por uma plataforma Web via navegador, com o fito de que os usuários a encontrem sem uma barreira inicial de instalação de um software específico, além de usufruir do alcance da Internet. A aplicação deve dispor de: interfaces práticas e intuitivas aos usuários, banco de dados, sistema de autenticação, gerenciamento de anúncios, integração com serviços de pagamento e mecanismos de pesquisa e filtragem de itens.

O sistema deverá oferecer funcionalidades voltadas tanto para compradores quanto para vendedores, permitindo que usuários publiquem anúncios contendo imagens, descrição, preço, categoria e estado de conservação das peças anunciadas. Os compradores deverão poder pesquisar produtos, adicionar itens ao carrinho, realizar compras e se comunicar diretamente com vendedores através do chat interno da plataforma.

A arquitetura da aplicação deverá priorizar escalabilidade, segurança e responsividade, permitindo o acesso em dispositivos desktop, tablets e smartphones.

\subsection{Definições \& abreviações}
\label{sec:Definicoes}

\begin{enumerate}[label=\alph*)]
    \item \textbf{Software Requirements Specification – SRS:} Especificação de Requisitos de Software.
    \item \textbf{Consumer-to-Consumer – C2C:} Modelo de E-Commerce entre consumidores.
    \item \textbf{Usuário:} Indivíduo que utiliza o sistema.
    \item \textbf{Sistema:} Plataforma de intermediação de brechós.
    \item \textbf{Plataforma:} Ambiente digital de uso dos usuários finais.
    \item \textbf{Item:} Peça de vestuário anunciada.
    \item \textbf{Anúncio:} Publicação de um item.
    \item \textbf{Carrinho:} Lista de itens adicionados para compra.
    \item \textbf{Chat:} Canal de comunicação entre usuários.
\end{enumerate}

\subsection{Referências}
\label{sec:Referencias}

IEEE Std 830-1998. IEEE Recommended Practice for Software Requirements Specifications. (1998)

Roger S. Pressman \& Bruce R. Maxim. Software Engineering: A Practitioner's Approach. 8th Edition. (2014)


\section{Descrição Geral}
\label{sec:DescricaoGeral}

\subsection{Perspectiva do produto}
\label{sec:PerspectivaProduto}

\subsubsection{Interfaces de Sistema}
\label{sec:InterfaceSistema}

O sistema de intermediação de brechós é um produto independente e totalmente contido. Ele não atua como um subsistema ou componente subordinado a um sistema maior. Todas as funcionalidades descritas neste documento representam a totalidade do ecossistema do produto.


\subsubsection{Interfaces de Hardware}
\label{sec:InterfaceHardware}

Não se aplica. O sistema não exige integração direta ou controle de baixo nível com nenhum equipamento ou periférico físico específico.

\subsubsection{Interfaces de Software}
\label{sec:InterfaceSoftware}

O sistema dependerá da integração com os seguintes softwares e serviços externos para o seu pleno funcionamento:

\begin{itemize}
    \item \textbf{Sistema Gerenciador de Banco de Dados (SGBD):}
    \begin{itemize}
        \item[ ] \textbf{Nome:} PostgreSQL
        \item[ ] \textbf{Versão:} 17.0 em ambiente local via Docker Compose; compatível com PostgreSQL gerenciado em produção.
        \item[ ] \textbf{Propósito:} Persistência estruturada e relacional dos dados do sistema, incluindo usuários, endereços, anúncios, fotos, carrinhos, compras, pedidos, entregas, pagamentos, reembolsos, conversas e notificações.
        \item[ ] \textbf{Interface:} Comunicação realizada pelo \textit{backend} Go com o driver pgx/v5, operando localmente ou em servidor gerenciado pela porta TCP padrão 5432.
    \end{itemize}
    
    \vspace{0.5em}
    
    \item \textbf{Gateway de Pagamentos:}
    \begin{itemize}
        \item[ ] \textbf{Nome:} Mercado Pago, com provedor simulado como contingência no ambiente acadêmico.
        \item[ ] \textbf{Versão:} API REST do Mercado Pago.
        \item[ ] \textbf{Propósito:} Processamento de cobranças, confirmação assíncrona por \textit{webhook}, pagamento simulado em ambiente sem credenciais e registro de reembolsos simulados. O \textit{split}/repasse financeiro real ao vendedor permanece fora do escopo implementado.
        \item[ ] \textbf{Interface:} Integração remota via protocolo HTTPS (porta 443), com mensagens em JSON. A autenticidade do \textit{webhook} é validada por assinatura configurada no ambiente.
    \end{itemize}
    
    \vspace{0.5em}

    \item \textbf{Serviço de Armazenamento de Arquivos:}
    \begin{itemize}
        \item[ ] \textbf{Nome:} Vercel Blob.
        \item[ ] \textbf{Versão:} API HTTP do Vercel Blob.
        \item[ ] \textbf{Propósito:} Armazenamento persistente e distribuição pública das fotos dos anúncios. O PostgreSQL armazena apenas URLs públicas e metadados das imagens.
        \item[ ] \textbf{Interface:} O \textit{backend} emite autorizações temporárias de upload para usuários autenticados; o navegador envia as imagens diretamente ao Blob store via HTTPS. O sistema aceita apenas URLs do host público configurado.
    \end{itemize}

    \vspace{0.5em}

    \item \textbf{Serviços de Frete e Endereço:}
    \begin{itemize}
        \item[ ] \textbf{Nome:} Melhor Envio para cotação de frete, com frete fixo como contingência; ViaCEP para consulta de CEP.
        \item[ ] \textbf{Propósito:} Cotar frete no checkout, validar endereços e preencher dados de endereço a partir do CEP informado.
        \item[ ] \textbf{Interface:} Integração remota via HTTPS. Na ausência de credenciais do Melhor Envio, o sistema utiliza cotação fixa para manter o fluxo de compra executável.
    \end{itemize}
\end{itemize}

\subsubsection{Interfaces de Comunicação}
\label{sec:InterfaceComunicacao}

A comunicação entre cliente e servidor deverá ocorrer através do protocolo HTTPS, garantindo criptografia e segurança na troca de informações. A implementação atual utiliza renderização no servidor (SSR) com HTMX para navegação, formulários, notificações e conversa entre comprador e vendedor. O chat é persistido no PostgreSQL e atualizado por requisições HTTP; não há servidor WebSocket dedicado nem envio SMTP implementado nesta versão.

\subsubsection{Restrições de Memória}
\label{sec:RestricoesMemoria}

O sistema deverá operar dentro de limites lógicos de memória primária e armazenamento para garantir a estabilidade da aplicação, evitar travamentos e otimizar os custos de infraestrutura:

\begin{itemize}
    \item \textbf{Limites de Upload (Armazenamento):} Para evitar a saturação do servidor e do banco de dados, o \textit{upload} de arquivos de mídia (como fotos para o cadastro de anúncios e avatares de perfil) deverá ser estritamente limitado a um tamanho máximo de 5MB por arquivo. O sistema permitirá um limite máximo de 5 imagens por anúncio.
    
    
    \item \textbf{Paginação de Dados em Memória:} Para evitar a sobrecarga de memória RAM tanto no servidor quanto no navegador do usuário, listagens extensas de dados nunca devem ser carregadas integralmente na memória. Componentes como o \textit{Feed} de Anúncios, Histórico de Pedidos e o Chat devem utilizar paginação estrita (ex: carregar pacotes de 20 a 30 registros por requisição).
\end{itemize}

\subsubsection{Operações}
\label{sec:Operacoes}

O sistema deverá permitir operações de cadastro, autenticação, criação de anúncios, busca de produtos, adição ao carrinho, finalização de compras, troca de mensagens entre usuários e avaliação de vendedores.

\subsubsection{Requisitos de Adaptação do Local}
\label{sec:RequisitosLocal}

O sistema deverá estar adaptado ao idioma português brasileiro, utilizando moeda em Real (R\$), formatos locais de data e compatibilidade com endereços nacionais.

\subsection{Funções do produto}
\label{sec:FuncoesProduto}

As principais funções do Sistema de Intermediação de Brechós incluem:

\begin{itemize}
    \item Cadastro e autenticação de usuários;
    \item Gerenciamento de perfil do usuário;
    \item Publicação, edição e exclusão de anúncios;
    \item Pesquisa e filtragem de produtos;
    \item Comunicação entre compradores e vendedores;
    \item Adição de produtos ao carrinho;
    \item Processamento de compras;
    \item Histórico de compras e vendas;
    \item Gerenciamento de pedidos;
    \item Upload de imagens para anúncios.
\end{itemize}

\subsection{Características do usuário}
\label{sec:CaracteristicasUsuario}

O sistema destina-se a um perfil de usuário que demonstre interesse em comprar e/ou vender peças de vestuário na plataforma. Um mesmo usuário pode atuar como comprador e vendedor de acordo com a sua preferência e, portanto, o sistema deve ser estruturado de forma a contemplar ambos casos de utilização. Espera-se que os usuários possuam um conhecimento básico a intermediário no uso de aplicações digitais, bem como um nível básico de familiaridade com plataformas de comércio eletrônico.

Nesse contexto, a interface do sistema deve ser prática e intuitiva o suficiente para favorecer a plena realização das principais operações por usuários, a exemplo de busca por categoria ou compra de itens.

\subsection{Restrições gerais}
\label{RestricoesGeral}
As restrições gerais do sistema são divididas sobre as seguintes categorias: Tecnológica, Operacional, Legal e de Integração. A tabela com a descrição das restrições e suas respectivas categorias encontra-se listada no Apêndice.

\subsection{Suposições e dependências}
\label{sec:Suposicoes}

O sistema pressupõe que os usuários possuam acesso estável à Internet e dispositivos compatíveis com navegação web. Além disso, determinadas funcionalidades dependem de serviços externos, como gateways de pagamento, serviços de hospedagem de imagens, e gateways de rastreio de entregas.

\section{Requisitos Específicos}
\label{RequisitosEspecificos}

\subsection{Requisitos funcionais}
\label{RequisitosFuncionais}

\begin{enumerate}[label=RF\arabic*:]

    \item O sistema deve permitir o cadastro de Contas mediante o fornecimento obrigatório de Nome Completo, CPF, E-mail, Senha e Endereço Completo.

    \item O sistema deve validar o CPF durante o cadastro.
    
    \item O sistema deve permitir autenticação de acesso via E-mail/CPF e Senha.

    \item O sistema deve permitir que os usuários criem anúncios.

    \item O sistema somente deve criar um anúncio mediante o fornecimento dos seguintes detalhes: Título, Descrição, Categoria, Tamanho, Cor, Estado de Conservação e Preço.

    \item O sistema deve exigir que, no mínimo, duas fotografias sejam vinculadas a um anúncio.

    \item O sistema deve permitir que, no máximo, cinco fotografias sejam vinculadas a um anúncio.

    \item O sistema deve permitir edição de anúncios ativos.

    \item O sistema deve permitir a exclusão de anúncios ativos. 

    \item O sistema deve permitir a pesquisa de anúncios por filtros de Categoria, Tamanho, Faixa de Preço ou de Estado de Conservação.

    \item O sistema deve permitir a pesquisa de anúncios por palavras-chave.

    \item O sistema deve permitir a comunicação em tempo real via chat de mensagens entre usuários.

    \item O sistema deve permitir a associação futura de dados financeiros do vendedor por meio de identificador opaco do provedor de pagamento, sem persistir dados bancários ou de cartão completos na aplicação.

    \item O sistema deve gerenciar um carrinho de compras por usuário.

    \item O sistema deve processar o pagamento do carrinho de compras em uma única transação financeira para o comprador, independentemente de os itens pertencerem a vendedores distintos.

    \item O sistema deve validar os status dos itens de um carrinho na hora da finalização da compra.

    \item O sistema deve permitir que os usuários gerenciem seus pedidos ativos, desde o acompanhamento do status do pedido de compra até o acompanhamento individual do status dos itens dentro do pedido.

    \item O sistema deve exibir o histórico de compras concluídas do usuário.

    \item O sistema deve exibir o histórico de vendas concluídas do usuário.

    \item O sistema deve registrar a elegibilidade de repasse ao vendedor somente após a confirmação de entrega do item. O repasse financeiro real ainda não está automatizado na implementação atual.

    \item O sistema deve permitir que o comprador confirme manualmente o recebimento do item adquirido. 

    \item O sistema deve informar o código de rastreio para o pedido após o pedido ser alterado para \textit{Enviado}.

    \item O sistema deve registrar o código de rastreio informado pelo vendedor e associá-lo ao pedido. A validação externa do rastreio e da compatibilidade entre CEPs permanece como evolução futura.

    \item O sistema deve realizar o reembolso para o comprador do valor do item específico  (incluindo a taxa de serviço) caso o item receba o status de \textit{Não Enviado}. 

    \item O sistema deve alterar automaticamente o status de um item de \textit{Não Enviado} para \textit{Suspenso} após o reembolso.

    \item O sistema deve bloquear o vendedor de continuar com suas atividades após, no máximo, três pedidos não enviados.

    \item O sistema deve permitir que o vendedor de um item altere o status desse item de \textit{Suspenso} para \textit{Disponível} mediante o pagamento de uma taxa proporcional ao valor da compra.

    \item O usuário deve ser notificado quando receber uma mensagem no chat, quando realizar uma venda ou quando o status do seu pedido mudar.

    \item O valor do frete deverá ser calculado individualmente para cada vendedor do carrinho (com base nos CEPs de origem e destino) e a soma desses fretes deverá compor o valor total pago pelo comprador no momento da compra.

     \item O sistema calculará automaticamente uma taxa de serviço percentual sobre o valor total da venda dos itens de um pedido, registrando o valor líquido previsto ao vendedor após a conclusão do pedido.

     \item O sistema deverá permitir que compradores avaliem (notas de 1 a 5 e comentários) os vendedores após a confirmação do recebimento do pedido.

\end{enumerate}

\subsection{Requisitos não-funcionais}
\label{RequisitosNaoFuncionais}

\begin{enumerate}[label=RNF\arabic*:]

    \item O sistema deverá apresentar tempo de carregamento inferior a 3 segundos na renderização inicial do feed de anúncios.

    \item O sistema deverá proteger a privacidade e os dados sensíveis dos usuários em conformidade com a Lei Geral de Proteção de Dados (LGPD).

    \item O sistema deverá salvar as senhas dos usuários com hash.

    \item O sistema deverá garantir plena compatibilidade de uso nas versões mais recentes dos navegadores Google Chrome, Mozilla Firefox, Safari e Microsoft Edge.

\end{enumerate}

\subsection{Requisitos de interface}
\label{RequisitosInterface}

\begin{enumerate}[label=RI\arabic*:]

    \item O sistema deverá apresentar alertas de confirmação para ações destrutivas (ex: exclusão de conta ou exclusão de anúncio).

    \item O sistema deve limitar o tamanho dos arquivos de upload (imagens) a até 5MB.

    \item Os formulários deverão validar dados e impedir o envio de requisições incompletas ao servidor.

\end{enumerate}

\subsection{Requisitos do domínio}
\label{RequisitosDominio}

\begin{enumerate}[label=RD\arabic*:]

    \item Um único usuário poderá exercer simultaneamente os papéis de Comprador e Vendedor na plataforma, utilizando a mesma conta.

    \item Todos os itens devem ser únicos, representando um único anúncio de peça de roupa.

    \item Uma vez excluído, um item deixa de existir para sempre no catálogo público.

    \item Um item só pode ser excluído ou alterado caso esteja no estado \textit{Disponível}.
    
    \item O ciclo de vida de um item deverá transitar estritamente pelos seguintes status: \textit{Disponível, Indisponível, Vendido, Enviado, Não Enviado, Suspenso e Recebido}.

    \item Cada anúncio deverá estar associado, de maneira exclusiva, a um único perfil de vendedor e um único item.

    \item Um usuário pode adicionar diversos itens ao seu carrinho de compras.

    \item A finalização do carrinho de compras gerará um ou mais pedidos de compra, de forma que cada pedido agrupe exclusivamente os itens pertencentes a um único vendedor.

     \item A adição de um item ao carrinho não garante a sua reserva física; a reserva do item no sistema deverá ocorrer exclusivamente no momento de geração do pedido de compra.

    \item O ciclo de vida de um pedido de compra deverá transitar estritamente pelos seguintes status: \textit{Criado, Cancelado, Aguardando Envio, Aguardando Entrega e Finalizado}.

    \item Um item inserido em um pedido deverá ser marcado temporariamente como \textit{Indisponível}, ficando oculto no catálogo público.

    \item Caso todos os itens pertencentes a um mesmo pedido de compra assumam o status de \textit{Suspenso}, o pedido correspondente deverá ter seu status automaticamente alterado para \textit{Cancelado}.

    \item Itens que assumirem o status de \textit{Suspenso} (devido à falha no envio) ficarão aptos a retornar ao status de \textit{Disponível} no catálogo exclusivamente após o pagamento da taxa pelo respectivo vendedor, independentemente do status geral do pedido de compra ao qual pertenciam.

    \item Após a confirmação de pagamento, os itens da compra deverão assumir o status de \textit{Vendido}, permanecendo ocultos do catálogo público. Simultaneamente, os pedidos de compra correspondentes deverão receber o status de Aguardando Envios. 

    \item Caso o vendedor de um item não altere o status do item para \textit{Enviado} no tempo mínimo de 7 dias, o item deve ser marcado como \textit{Não Enviado}.

    \item Itens marcados como \textit{Suspenso} ficam ocultos do catálogo público. 

    \item Quando todos os itens de um pedido forem marcados como \textit{Enviado} pelos seus respectivos vendedores ou \textit{Suspenso} pelo sistema, o pedido deve receber o status de \textit{Aguardando Entregas}.

    \item O pedido de compra recebe o status de \textit{Finalizado} quando todos os itens do pedido são marcados como \textit{Recebido} pelo comprador ou \textit{Suspenso} pelo sistema.      

    \item Caso o comprador não confirme o recebimento manualmente em até 15 dias corridos após o pedido ser marcado como \textit{Enviado}, o sistema poderá alterar o status do item para \textit{Recebido} automaticamente e registrar a elegibilidade de repasse ao vendedor.

\end{enumerate}


\clearpage
\section{Apêndice}
\label{Apendice}

\begin{longtable}{|l|p{10cm}|}

\caption{Restrições Gerais do Sistema} \\

\hline
\cellcolor{gray!30} \textbf{Categoria} &
\cellcolor{gray!30} \textbf{Descrição da Restrição} \\
\hline
\endfirsthead

\hline
\cellcolor{gray!30} \textbf{Categoria} &
\cellcolor{gray!30} \textbf{Descrição da Restrição} \\
\hline
\endhead

Tecnológica & Ofertar compatibilidade com navegadores modernos: Chrome (ver. 147.0.7727.138), Firefox (ver. 150.0.2), Safari (ver. 26) e Microsoft Edge (ver. 148.0.3967.54). \\
\hline

Tecnológica & Limitação do tamanho máximo de imagens enviadas pelos usuários e uso de técnicas de compressão (Restrição de Memória). \\
\hline

Tecnológica & O acesso à plataforma dependerá de conexão com a Internet. \\
\hline

Operacional & O sistema deve estar sujeito ao acesso contínuo dos usuários, exceto em períodos de manutenção programada. \\
\hline

Operacional & O sistema deverá suportar o acesso de múltiplicos usuários simultaneamente. \\
\hline

Operacional & A plataforma é isenta de responsabilidade sobre a logística de envios, qualidade física e higienização das peças transacionadas. \\
\hline

Operacional & Adaptação somente ao idioma Português Brasileiro (PT-BR), moeda Real (R\$) e formatos locais de data e endereço. \\
\hline

Legal & Conformidade com a LGPD (Lei nº 13.709/2018) para armazenamento e tratamento de dados pessoais e sensíveis dos usuários. \\
\hline

Legal & Conformidade com o Código de Defesa do Consumidor (Lei nº 8.078/1990) para transparência e responsabilidade com os usuários. \\
\hline

Legal & Conformidade com o Marco Civil da Internet (Lei nº 12.965/2014) para proteção dos direitos dos usuários no sistema. \\
\hline

Legal & Presença obrigatória de Termos de Uso e Políticas de Privacidade na plataforma. \\
\hline

Organizacional & A documentação do sistema deve estar atualizada e relacionada com o software em produção. \\
\hline

Organizacional & As atividades do projeto devem seguir o cronograma estabelecido. \\
\hline

Organizacional & As equipe de desenvolvimento deve aplicar padronização sobre o código-fonte. \\
\hline

Organizacional & O projeto deve seguir as diretrizes e princípios da Engenharia de Software. \\
\hline

Integração & O sistema dependerá da disponibilidade de serviços externos (Gateways de pagamento e hospedagem em nuvem). \\
\hline

Integração & As integrações externas devem ser compatíveis com as tecnologias Web utilizadas para o sistema. \\
\hline

Integração & O processamento e validação de transações financeiras ocorrerão fora do sistema, via API de terceiros (Gateways de pagamento). \\
\hline

Integração & Dependência de serviços externos (hospedagem em nuvem) para armazenamento de imagens enviadas nos anúncios. \\
\hline

\end{longtable}



% ----------------------------------------------------------
% Índice
% ----------------------------------------------------------
\clearpage
\section{Índice}
\label{Indice}

Introdução\dotfill \pageref{sec:Introducao}

\hspace{1cm}Propósito do documento\dotfill \pageref{sec:Proposito}

\hspace{1cm}Escopo do produto\dotfill \pageref{sec:Escopo}

\hspace{2cm}Produtos de software utilizados\dotfill \pageref{sec:Produtos}

\hspace{2cm}O que o sistema realiza (e o que não realiza)\dotfill \pageref{sec:Realiza}

\hspace{2cm}Descrição da aplicação de software proposta\dotfill \pageref{sec:DescricaoAplicacao}

\hspace{1cm}Definições \& abreviações\dotfill \pageref{sec:Definicoes}

\hspace{1cm}Referências\dotfill \pageref{sec:Referencias}

Descrição Geral\dotfill \pageref{sec:DescricaoGeral}

\hspace{1cm}Perspectiva do produto\dotfill \pageref{sec:PerspectivaProduto}

\hspace{2cm}Interfaces de Sistema\dotfill \pageref{sec:InterfaceSistema}

\hspace{2cm}Interfaces de Hardware\dotfill \pageref{sec:InterfaceHardware}

\hspace{2cm}Interfaces de Software\dotfill \pageref{sec:InterfaceSoftware}

\hspace{2cm}Interfaces de Comunicação\dotfill \pageref{sec:InterfaceComunicacao}

\hspace{2cm}Restrições de Memória\dotfill \pageref{sec:RestricoesMemoria}

\hspace{2cm}Operações\dotfill \pageref{sec:Operacoes}

\hspace{2cm}Requisitos de Adaptação do Local\dotfill \pageref{sec:RequisitosLocal}

\hspace{1cm}Funções do produto\dotfill \pageref{sec:FuncoesProduto}

\hspace{1cm}Características do usuário\dotfill \pageref{sec:CaracteristicasUsuario}

\hspace{1cm}Restrições gerais\dotfill \pageref{RestricoesGeral}

\hspace{1cm}Suposições e dependências\dotfill \pageref{sec:Suposicoes}

Requisitos Específicos\dotfill \pageref{RequisitosEspecificos}

\hspace{1cm}Requisitos funcionais\dotfill \pageref{RequisitosFuncionais}

\hspace{1cm}Requisitos não-funcionais\dotfill \pageref{RequisitosNaoFuncionais}

\hspace{1cm}Requisitos de interface\dotfill \pageref{RequisitosInterface}

\hspace{1cm}Requisitos do domínio\dotfill \pageref{RequisitosDominio}

Apêndice\dotfill \pageref{Apendice}

Índice\dotfill \pageref{Indice}

\end{document}
